\documentclass[10pt,a4paper]{amsart}
\usepackage[margin=19mm]{geometry}
\usepackage{amsmath,amssymb,mathtools}
\usepackage[scr=rsfs,cal=euler]{mathalfa}
\usepackage{xfrac}
\usepackage[unicode,hidelinks,bookmarks=false]{hyperref}
\newcommand{\ie}{i.e.\ }
\newcommand{\cf}{cf.\ }
\newcommand{\resp}{respectively\ }
\newcommand{\Subsection}{\subsection}
\providecommand{\nobreakdash}{\nobreak\hbox{-}}
\setlength{\emergencystretch}{2em}
\multlinegap0pt
\usepackage{bm}
\usepackage{bbm}
\usepackage{dsfont}
\usepackage{xspace}
\usepackage{cancel}
\usepackage{mathtools}
\usepackage{amsthm}
\makeatletter
\@ifundefined{corollary}{\newtheorem{corollary}{Corollary}}{}
\makeatother
\makeatletter
\expandafter\ifx\csname oneshot\endcsname\relax
\newenvironment{oneshot}[1]{\@begintheorem{#1}{\unskip}}{\@endtheorem}
\fi
\makeatother
\title[Composition and Coherence: The Syntax of Operator Networks]{Composition and Coherence: The Syntax of Operator Networks}
\author{Shih-Yu Chang}
\date{Source: arXiv:2511.13706v1 (CC BY 4.0)}
\begin{document}
\maketitle

\noindent\emph{Source and licence.} Composition and Coherence: The Syntax of Operator Networks. Version arXiv:2511.13706v1 (CC BY 4.0), distributed under the Creative Commons Attribution 4.0 International licence, as recorded in the arXiv licence metadata for this exact version and shown on the abstract page https://arxiv.org/abs/2511.13706v1. Full rights notice: licenses/arxiv-2511.13706-CC-BY-4.0.txt.\par\smallskip

\noindent\textit{Editorial scope.} Counted unit: the Corollary whose source label is \texttt{cor:differentiable-feedback} in the article (source0.tex lines 1193--1201, source label \texttt{cor:differentiable-feedback}), delivered together with the article's own complete original proof of it (source0.tex lines 1203--1297, one \texttt{proof} environment of 95 lines). No mathematics is written, repaired or re-derived: statement and proof are the article's own text line for line. Required context retained uncounted: cor:feedback-stability (lines 1017--1031); cor:control-continuity (lines 1112--1118). Every definition, display and label of those lines is retained unchanged. Omitted from this package and not redistributed: 1--1016, 1032--1111, 1119--1192, 1202--1202, 1298--2206. Nothing inside a retained excerpt is omitted or altered: every retained body is delivered exactly as the article's lines are written, so no omission receipt of any kind accompanies this package. Citations: the retained excerpts carry no citation command at all, so no bibliography entry of the article is transcribed and no reference is redistributed. Reference adaptation: none was needed and none was made; every reference inside the delivered unit resolves inside it, so no unresolved reference or citation is produced. Printed slips kept literal: no printed word, symbol, hypothesis or number of the counted statement or of its proof was altered. Layout and class: the article's own class is \texttt{article} with packages including geometry, times, expl3, cite, xcolor, multirow, stackengine, hhline, lipsum, titlesec, wrapfig, enumerate, epsfig, tikz-cd; the wrapper is the lane's pinned \texttt{amsart} editorial wrapper at 10pt on A4, which restates the theorem-like environments the retained text needs and adds the article's own macro definitions that the retained text needs (none), with extra wrapper packages (bm, bbm, dsfont, xspace, cancel, mathtools). The delivered numbering is the unit's own. Assets: no retained span contains a figure environment, a graphics-inclusion command, a PDF-page inclusion, a table or a TikZ picture; no image file is copied into this package, the package declares no asset dependency and no image file is redistributed with it. Licence: version arXiv:2511.13706v1 of this article is distributed under the Creative Commons Attribution 4.0 International licence, as recorded in the arXiv licence metadata for this exact version and shown on the abstract page https://arxiv.org/abs/2511.13706v1. A full rights notice is delivered at licenses/arxiv-2511.13706-CC-BY-4.0.txt. Characters: every retained excerpt is pure ASCII, so the delivered engine, XeLaTeX with its default Latin Modern text font, reports no missing character.
\begin{corollary}[Contractive Fixed-Point Theorem (Feedback Stability)]\label{cor:feedback-stability}
Let \( T: \rho(\vec{H}') \otimes X \to X \) be a bounded multilinear map in \(\mathbf{HilbMult}\), and consider the feedback equation for each \( p \in \rho(\vec{H}') \):
\[
\xi = T(p \otimes \xi) + b(p)
\]
where \( b: \rho(\vec{H}') \to X \) is a bounded multilinear map. If for each \( p \in \rho(\vec{H}') \) the mapping \( \Phi_p: \xi \mapsto T(p \otimes \xi) + b(p) \) is a contraction on \( X \) with uniform contraction constant \( \kappa < 1 \), then:

\begin{enumerate}
    \item The feedback operator \( \rho(\mathcal{F}_{i,j})(T,b): \rho(\vec{H}') \to X \) is well-posed, yielding a unique solution \( \xi(p) \) for each input \( p \)
    \item The resulting map \( p \mapsto \xi(p) \) is multilinear and bounded with norm estimate:
    \[
    \|\rho(\mathcal{F}_{i,j})(T,b)\| \le \frac{\|b\|}{1 - \kappa}
    \]
\end{enumerate}
\end{corollary}

\begin{corollary}[Continuity of Control Action]\label{cor:control-continuity}
Assume the control injection maps \( T(H_i, -): \mathcal{U} \to \mathcal{B}(H_i, H_i') \) are continuous in the operator norm topology, and the monoid multiplication \( \star: \mathcal{U} \times \mathcal{U} \to \mathcal{U} \) is continuous. Then the induced semantic control map
\[
\Phi: \mathcal{U} \to \mathcal{B}(\rho(\vec{H}), \rho(\vec{H}')), \quad \Phi(u) = \rho(\mathrm{Ctrl}_u)
\]
is continuous, endowing the control transformations with the structure of a topological semigroup action. Furthermore, if \( \mathcal{U} \) is a Lie group with smooth multiplication and the maps \( T(H_i, -) \) are smooth, then \( \Phi \) is smooth, yielding a Lie group action on the operator space.
\end{corollary}

\begin{corollary}[Differentiable Feedback Law]\label{cor:differentiable-feedback}
Suppose the control injection maps \( T_i: H_i \otimes \mathcal{U} \to H_i' \) are Fr\'echet differentiable with respect to \( u \in \mathcal{U} \), and the feedback context satisfies the well-posedness conditions of Corollary~\ref{cor:feedback-stability}. Then the combined semantic control-and-feedback map
\[
\Psi: \rho(\vec{H}) \times \mathcal{U} \longrightarrow X, 
\qquad 
\Psi(\psi, u) = \rho(\mathcal{F}_{i,j})(\rho(\mathrm{Ctrl}_u))(\psi)
\]
is Fr\'echet differentiable. Consequently, the family of maps \( \{\Psi_u: \rho(\vec{H}) \to X\}_{u \in \mathcal{U}} \) varies smoothly with the control parameter.
\end{corollary}

\begin{proof}
We analyze the differentiability by decomposing \(\Psi\) and applying the implicit function theorem.

\medskip
\noindent\textbf{Step 1: Correct Type Analysis and Decomposition}
The composition has the following structure:
\[
\rho(\vec{H}) \times \mathcal{U} 
\xrightarrow{\rho(\mathrm{Ctrl})} 
\rho(\vec{H}')
\xrightarrow{\rho(\mathcal{F}_{i,j})(\rho(\mathrm{Ctrl}_u))} 
X,
\]
where we define the parametrized feedback operator. More precisely, define the family of maps:
\[
T_u: \rho(\vec{H}') \otimes X \to X, \quad T_u(p \otimes \xi) = \rho(\mathrm{Ctrl}_u)(p \otimes \xi),
\]
so that \(\Psi(\psi, u) = \rho(\mathcal{F}_{i,j})(T_u)(\psi)\).

\medskip
\noindent\textbf{Step 2: Differentiability of the Parametrized Map \(T_u\)}
From Corollary~\ref{cor:control-continuity}, the control map is Fr\'echet differentiable in \(u\). Specifically, the map
\[
u \mapsto \rho(\mathrm{Ctrl}_u) \in \mathcal{B}(\rho(\vec{H}) \otimes \mathcal{U}, \rho(\vec{H}'))
\]
is Fr\'echet differentiable by the tensor product construction and the differentiability of each component \(T_i\).

Therefore, the parametrized map
\[
F: \mathcal{U} \to \mathcal{B}(\rho(\vec{H}') \otimes X, X), \quad F(u) = T_u
\]
is Fr\'echet differentiable.

\medskip
\noindent\textbf{Step 3: Differentiability of the Feedback Solution Map}
Define the solution map:
\[
S: \mathcal{B}(\rho(\vec{H}') \otimes X, X) \to \mathcal{B}(\rho(\vec{H}'), X), \quad S(T) = \rho(\mathcal{F}_{i,j})(T),
\]
which assigns to each \(T\) the unique fixed point \(\xi\) satisfying \(\xi(p) = T(p \otimes \xi(p))\) for all \(p \in \rho(\vec{H}')\).

Consider the implicit function:
\[
G: \mathcal{B}(\rho(\vec{H}') \otimes X, X) \times \mathcal{B}(\rho(\vec{H}'), X) \to \mathcal{B}(\rho(\vec{H}'), X)
\]
defined by
\[
G(T, \xi)(p) = \xi(p) - T(p \otimes \xi(p)).
\]

At a solution \((T_0, \xi_0)\) with \(G(T_0, \xi_0) = 0\), compute the Fr\'echet derivative with respect to \(\xi\):
\[
D_\xi G(T_0, \xi_0)[\eta](p) = \eta(p) - D_2T_0(p \otimes \xi_0(p))[\eta(p)],
\]
where \(D_2T_0\) denotes the derivative of \(T_0\) with respect to its second argument.

Under the well-posedness assumption from Corollary~\ref{cor:feedback-stability}, the operator
\[
L(\eta)(p) = D_2T_0(p \otimes \xi_0(p))[\eta(p)]
\]
satisfies \(\|L\| \le \kappa < 1\). Therefore, \(D_\xi G(T_0, \xi_0) = I - L\) is invertible with bounded inverse \(\|(I - L)^{-1}\| \le \frac{1}{1 - \kappa}\).

By the implicit function theorem in Banach spaces, there exists a neighborhood \(U\) of \(T_0\) and a unique Fr\'echet differentiable function
\[
S: U \to \mathcal{B}(\rho(\vec{H}'), X)
\]
such that \(G(T, S(T)) = 0\) for all \(T \in U\). This function is precisely the feedback solution map \(S(T) = \rho(\mathcal{F}_{i,j})(T)\).

\medskip
\noindent\textbf{Step 4: Differentiability of the Composition}
We now have:
\begin{itemize}
    \item \(F: \mathcal{U} \to \mathcal{B}(\rho(\vec{H}') \otimes X, X)\) is Fr\'echet differentiable (Step 2)
    \item \(S: \mathcal{B}(\rho(\vec{H}') \otimes X, X) \to \mathcal{B}(\rho(\vec{H}'), X)\) is Fr\'echet differentiable (Step 3)
\end{itemize}
Therefore, their composition \(S \circ F: \mathcal{U} \to \mathcal{B}(\rho(\vec{H}'), X)\) is Fr\'echet differentiable.

Finally, for fixed \(\psi \in \rho(\vec{H})\), the evaluation map
\[
\mathrm{ev}_\psi: \mathcal{B}(\rho(\vec{H}'), X) \to X, \quad \mathrm{ev}_\psi(f) = f(\psi)
\]
is bounded linear, hence smooth. Therefore, the map
\[
\Psi(\psi, u) = \mathrm{ev}_\psi(S(F(u))) = \rho(\mathcal{F}_{i,j})(\rho(\mathrm{Ctrl}_u))(\psi)
\]
is Fr\'echet differentiable jointly in \((\psi, u)\).

\medskip
\noindent\textbf{Step 5: Smooth Parameter Dependence}
For each \(u \in \mathcal{U}\), define \(\Psi_u: \rho(\vec{H}) \to X\) by \(\Psi_u(\psi) = \Psi(\psi, u)\). The joint Fr\'echet differentiability of \(\Psi\) implies that the map
\[
u \mapsto \Psi_u \in \mathcal{B}(\rho(\vec{H}), X)
\]
is Fr\'echet differentiable, establishing smooth dependence on the control parameter.
\end{proof}

\end{document}
